\section{引言：为什么要读这篇文章}

Anthropic 是 Claude 系列大语言模型的开发商，在过去一年中与数十个行业团队合作构建基于 LLM 的 Agent 系统。这篇文章凝聚了他们的第一手实践经验，核心发现是：\textbf{最成功的实现往往不依赖复杂框架，而是采用简单、可组合的设计模式}。

本文的独特价值在于：它不是一篇学术论文，而是来自工业界一线的「设计模式手册」。文章系统性地将 Agentic 系统分解为 \textbf{Workflow}（工作流）和 \textbf{Agent}（自主代理）两大类，详细介绍了五种 Workflow 模式和 Agent 的设计原则，并在附录中给出了实践案例和工具设计的最佳实践。

\begin{importantbox}{本文核心论点}
构建 LLM 应用的成功之道不在于打造最复杂的系统，而在于为需求打造\textbf{最合适}的系统。应从简单 prompt 出发，通过系统评估优化，仅在简单方案不足时才引入多步骤 agentic 系统。
\end{importantbox}

\subsection{本章小结}
本文来自 Anthropic 的工程实践总结，核心主张是「从简到繁」地构建 agentic 系统。理解这篇文章有助于在实际开发中做出正确的架构决策。


